\section{Completing sampled prefixes and recovering the budget}
\label{comp:section}

This section turns the sampled balls into a solution with the original
number of facilities.  Two kinds of removal destinations are treated
separately: destinations represented by proxies or balls, and destinations
at bad comparison facilities for which no ball has been sampled.  The
separation allows the latter destinations to use their actual distances
without paying for a ball around each of them.  We first give the
submodular optimization argument, then the completion construction, and
finally its application to the endpoint of Theorem~\ref{sam:successful-prefix}.
The construction adapts the guessed-ball savings formulation of
\cite[Section~2.3]{CGKLL2019} and the joint addition/removal surrogate
of \cite[Section~4.6]{CGLSS2026}.  The blue destination optimization
below separates centers without sampled balls from the remaining
submodular choices.  The finite maximization proof is a specialization
of continuous greedy and partition rounding
\cite[Theorem~1.1 and Section~3.3]{CCPV2011}.  We prove the complete
optimization, realization and recovery statements needed here.

\subsection{A finite partition-constrained maximization procedure}

For a finite set $E$, a function $g:2^E\to\mathbb R_{\ge0}$ is
\emph{normalized} if $g(\varnothing)=0$, \emph{monotone} if
$g(U)\le g(V)$ whenever $U\subseteq V$, and \emph{submodular} if
\[
 g(U\cup\{e\})-g(U)\ge g(V\cup\{e\})-g(V)
 \qquad(U\subseteq V,\ e\notin V).
\]
We write $\Delta_e g(U)=g(U\cup\{e\})-g(U)$, with value zero when
$e\in U$.  A partition constraint permits at most one element from each
part of a specified partition $E=E_1\sqcup\cdots\sqcup E_a$.

\begin{lemma}[Categorical continuous greedy with explicit error]
\label{comp:partition-max}
Let $g$ be a normalized, monotone, submodular function given by an exact
value oracle, and let $E$ have a partition into $a$ parts.  For every
$0<\varepsilon<1-1/e$, a randomized algorithm uses
$\operatorname{poly}(|E|,a,1/\varepsilon)$ oracle calls and returns a
feasible $X$ satisfying
\[
 \mathbb E[g(X)]\ge(1-1/e-\varepsilon)
       \max\{g(Y): |Y\cap E_j|\le1\text{ for all }j\}.
\]
Every random draw can be made using a bounded polynomial number of fair
bits.  If the parts are nonempty, filling unoccupied parts preserves the
guarantee.  In particular, one obtains the factor $0.62$; more generally,
for every fixed rational $\theta>1/e$ with $\theta<1$, one obtains
expected gain at least $1-\theta$ times the maximum.
\end{lemma}

\begin{proof}
Empty parts may be discarded.  Write $q=|E|$, let $g_{\max}$ denote the
maximum in the statement, and compute
$M_1=\max_{e\in E}g(\{e\})$ using $q$ oracle calls.  If $E$ is empty or
$M_1=0$, submodularity and telescoping give $g(Y)\le\sum_{e\in Y}
g(\{e\})=0$ for every $Y$, so the conclusion is immediate.  Otherwise
$a\ge1$ and $0<M_1\le g_{\max}$.

For vectors $x\ge0$ with $\sum_{e\in E_j}x_e\le1$, independently in
each part select $e$ with probability $x_e$, and select nothing with the
remaining probability.  Let $R_x$ be the resulting feasible random set
and $G(x)=\mathbb E[g(R_x)]$.  Multilinearity across the parts gives
\begin{equation}
 \partial_eG(x)=\mathbb E\bigl[
   \Delta_e g(R_x\setminus E_j)\bigr],\qquad e\in E_j.
 \label{comp:categorical-derivative}
\end{equation}
Every integrand is between zero and $M_1$.  For a fixed feasible
maximizer $O_g$, monotonicity and submodularity give, for every $R_x$,
\[
 g_{\max}-g(R_x)
 \le\sum_{e\in O_g}\Delta_e g(R_x)
 \le\sum_{e\in O_g}\Delta_e g(R_x\setminus E_{j(e)}).
\]
The last inequality also holds when $e\in R_x$, since its left summand
is then zero.  Thus
\begin{equation}
 \sum_{j=1}^a\max_{e\in E_j}\partial_eG(x)
       \ge g_{\max}-G(x).
 \label{comp:derivative-direction}
\end{equation}

Choose a dyadic rational $\tau$ of bit length $O(\log(2/\varepsilon))$ with
$\min\{\varepsilon/8,1/16\}\le\tau\le\min\{\varepsilon/4,1/8\}$ and
let $T$ be the smallest power of two at least $a^2/\tau$.  Starting at
$x^{(0)}=0$, perform $T$ steps.  At step $t$, estimate every derivative
in Equation~\eqref{comp:categorical-derivative}, choose a largest
estimated derivative in each part, and add $1/T$ to each chosen
coordinate.  Let $v^{(t)}$ be the zero--one vector of chosen coordinates.
At time $t$ every part has total mass $t/T$, so all iterates and all
segments $x^{(t)}+uv^{(t)}$, $0\le u\le1/T$, are feasible.

There are $Q_0=Tq$ derivative estimates.  For each estimate take the
mean of
\[
 L=\left\lceil\frac{Q_0a^2}{4\tau^3}\right\rceil
\]
independent samples of its marginal.  Each sample has variance at most
$M_1^2/4$.  Chebyshev's inequality shows that the probability of error
greater than $\tau M_1/a$ is at most $\tau/Q_0$.  This assertion holds
conditionally on every preceding choice, because fresh samples are
used at the current iterate.  The union bound therefore shows that all
estimates are accurate with probability at least $1-\tau$, even though
the iterates are adaptive.  On this event, the chosen direction obeys
\begin{equation}
 \nabla G(x^{(t)})\cdot v^{(t)}
 \ge g_{\max}-G(x^{(t)})-2\tau M_1.
 \label{comp:estimated-direction}
\end{equation}

For completeness, the finite step has a controlled integration error.
Couple a categorical draw at $x$ and at $x+uv$ by changing a part's
outcome only when the additional mass $u$ assigned to its selected
coordinate is used.  Such a change occurs with probability $u$ in each
part.  For a derivative with its own part omitted, the probability that
any other part changes is at most $au$.  Since the marginal in
Equation~\eqref{comp:categorical-derivative} lies in $[0,M_1]$, its
expectation changes by at most $auM_1$.  Summing over the $a$ selected
coordinates and integrating for $0\le u\le1/T$ gives, with the slightly
looser bound convenient below,
\[
 G(x^{(t+1)})-G(x^{(t)})
 \ge\frac{g_{\max}-G(x^{(t)})-2\tau M_1}{T}
       -\frac{a^2M_1}{T^2}
 \ge\frac{g_{\max}-G(x^{(t)})-3\tau M_1}{T}.
\]
Iterating from $G(0)=0$ yields
\[
 G(x^{(T)})\ge
 \bigl(1-(1-1/T)^T\bigr)(g_{\max}-3\tau M_1)
 \ge(1-1/e-3\tau)g_{\max}.
\]
Here $g_{\max}-3\tau M_1\ge0$.  Sample the final categorical set
independently after the estimates.  On estimation failure its gain is
still nonnegative.  Consequently its unconditional expected gain is at
least
\[
 (1-\tau)(1-1/e-3\tau)g_{\max}
 \ge(1-1/e-4\tau)g_{\max}
 \ge(1-1/e-\varepsilon)g_{\max}.
\]
All categorical probabilities in this algorithm have denominator $T$;
one uniform integer in $\{0,\ldots,T-1\}$ therefore samples each part
exactly using $\log_2T$ fair bits.  The oracle count and the number of
such draws are polynomial in the stated parameters.  In the present
application oracle values are rational numbers of polynomial bit
length, so the arithmetic also takes polynomial time.  Monotonicity
justifies filling parts.  Finally $1-1/e>0.6321$, so an error smaller
than $0.0121$ gives $0.62$.  For the last assertion choose a rational
$\varepsilon$ with $0<\varepsilon\le\theta-1/e$.
\end{proof}

\subsection{The blue optimization and the red surrogate}

Fix $|S|=|O|=h\ge1$ and an injective proxy map $i\mapsto\phi_i\in S$
for the good members of $O$.  Let $m>0$ be the number of bad members and
let $S_0$ be the $m$ members of $S$ not used as proxies.  Consider a
prefix with $Q\subseteq S_0$ and one ball
\[
 \mathcal F_i=\{x\in F:d(w_i,x)\le R_i\}
\]
for each $i$ in a subset $A$ of the bad centers.  Assume $i\in\mathcal F_i$,
put $r_i=d(w_i,i)$, and let $B$ be the remaining bad centers.  The
algorithm sees the balls but not their indices in $O$.  As in the
sampling section, $\delta_i$ is the leaf at distance $R_i$ from $w_i$,
$\Lambda=\{\delta_i:i\in A\}$, and $\delta_\infty$ is the permanent
dummy.  Thus
\[
 H=(S\setminus Q)\cup\Lambda\cup\{\delta_\infty\}.
\]
Keep the fixed nearest assignments in $H$.  For a real $s\in S\setminus Q$,
write $J_s$ for its assigned clients and $w'_s=|J_s|$.

Set $a=|A|$, $c=m-a$, and $r=m-|Q|$.
The eventual count is illustrated in Figure~\ref{fig:completion-budget}:
the completion removes $r$ more original facilities and adds at most
$a+c=m$ real destinations.  We now optimize their choices.
\begin{figure}[tbp]
\centering
\begin{tikzpicture}[
  box/.style={draw=MidnightBlue,rounded corners=2pt,align=center,
    text width=4.05cm,minimum height=2.05cm,inner sep=5pt,font=\small},
  arrow/.style={-{Stealth[length=2mm]},semithick},
  node distance=7mm]
\node[box] (anchor) {\textbf{Real anchor}\\[4pt]
  $S\subseteq F$\\$|S|=h$};
\node[box,right=of anchor] (reference) {\textbf{Sampled reference}\\[4pt]
  $h-|Q|$ real facilities\\
  $a$ leaves and one\\permanent dummy};
\node[box,right=of reference] (output) {\textbf{Real completion}\\[4pt]
  remove $r=m-|Q|$ more\\
  add at most $a+c=m$\\
  keep at most $h$ facilities};
\draw[arrow] (anchor)--(reference);
\draw[arrow] (reference)--(output);
\node[below=8mm of reference,align=center,text width=14.8cm,font=\small]
  {$\displaystyle h-|Q|-r+a+c=h$\quad before merging repeated choices.\\[3pt]
   Ball destinations: at most $a$; destinations without a ball: at most $c=m-a$.};
\end{tikzpicture}
\caption{The facility count in completion.  The reference contains virtual
leaves, but the output contains only members of $F$.  A chosen facility
in a ball replaces its leaf; a nonempty real solution dominates the
permanent dummy.  The removals are distinct, while additions are set
unions, so overlaps and repeated destinations only reduce the count.
The figure describes accounting, not a Euclidean embedding or a
disjointness assumption on facility balls.}
\label{fig:completion-budget}
\end{figure}

Color every member of
$S\setminus Q$ green, red, or blue independently.  Green has probability
$1/2$ and each of red and blue has probability $1/4$, so two fair bits
suffice.  Enumerate all nonnegative counts
$b_{\mathrm{red}}+b_{\mathrm{blue}}=r$.  A count with too few candidates
of a required color is infeasible and is discarded.

The following blue-label construction adapts the color-coding method
of Alon--Yuster--Zwick~\cite[Section~3, Lemma~3.1]{AlonYusterZwick1995}.
The weighted block recurrence and its routing-cost analysis are given here.
For a fixed count $b=b_{\mathrm{blue}}>0$, assign each blue candidate an
independent label in $[b]$.  It suffices to use a dyadic distribution
whose probability for every label is at least $1/(2b)$: sample a uniform
integer in $\{0,\ldots,2^{\lceil\log_2 b\rceil}-1\}$ and use one plus
its remainder modulo $b$.  For a nonempty $U\subseteq[b]$, define
\begin{equation}
 C(U)=\min_{x\in F}\sum_{\ell\in U}
       \min_{\substack{s\text{ blue}\\\operatorname{label}(s)=\ell}}
              w'_s d(s,x).
 \label{comp:blue-block}
\end{equation}
A missing label makes this value infinite.  For every block also store
minimizing witnesses $x$ and one removal candidate of each label.
The subset dynamic program
\[
 D(0,\varnothing)=0,\qquad
 D(t,U)=\min_{\varnothing\ne V\subseteq U}
             \bigl(C(V)+D(t-1,U\setminus V)\bigr)
\]
has all other zero-block states infinite; a minimum over an empty
collection is also infinite.  Take
$K=\min_{1\le t\le c}D(t,[b])$, and discard the branch if $K=\infty$.
It selects exactly one blue removal per label and at most $c$ destination
facilities.  If $b>0$ and $c=0$ the branch is infeasible.  If $b=0$, set
$K=0$ and select no blue removal or destination, including when $c=0$.

The block costs require at most $2^b\operatorname{poly}(N)$ work, and
the dynamic program requires $O(c3^b)$ additional arithmetic operations.
It is exact for this labeled optimization: every feasible routing
partitions the labels by destination and hence has cost at least a
corresponding sum of block minima; conversely each finite partition has
the stored witnesses realizing its value.  Distinct labels select
distinct removals, because a candidate has only one label.  Different
blocks may use the same destination, which only reduces the number of
distinct facilities opened.

Let $E=\bigsqcup_{i\in A}E_i$ consist of disjoint copies of the
facilities in the balls, with $E_i=\{(i,x):x\in\mathcal F_i\}$.
Distances to a subset of $E$ mean distances to its underlying facilities.
Ball intersections cause no identification between these copies.
For a feasible $X\subseteq E$, put
\[
 Z=\{s:s\text{ green}\}\cup\Lambda\cup\{\delta_\infty\}
\]
and define the finite surrogate
\begin{equation}
 \Psi(X)=K+\sum_{p\in D}\min\{d(p,H),d(p,X)\}
  +\operatorname{small}_{b_{\mathrm{red}}}
       \{w'_s d(s,Z\cup X):s\text{ red}\}.
 \label{comp:surrogate}
\end{equation}
Here $d(p,\varnothing)=\infty$, and $\operatorname{small}_q$ denotes
the sum of the $q$ smallest entries, with $\operatorname{small}_0=0$.

\begin{lemma}[Surrogate realization and submodularity]
\label{comp:surrogate-properties}
For every finite branch and feasible $X$, one can construct a nonempty
real solution $U_X\subseteq F$ with $|U_X|\le h$ and
$\operatorname{cost}(U_X)\le\Psi(X)$.  Moreover
$g(X)=\Psi(\varnothing)-\Psi(X)$ is normalized, nonnegative, monotone,
and submodular on all subsets of $E$.
\end{lemma}

\begin{proof}
Select the $b_{\mathrm{red}}$ red removals attaining the last term of
Equation~\eqref{comp:surrogate}, and use the blue removals and destinations
stored by the dynamic program.  Fill every unused ball with an arbitrary
facility in that ball.  Let $X_{\mathrm{fill}}$ be the underlying filled
facility set, $Y_{\mathrm{blue}}$ the blue destinations, and
$R_{\mathrm{red}},R_{\mathrm{blue}}$ the two removal sets.  Form
\[
 U'=\bigl(S\setminus(Q\cup R_{\mathrm{red}}\cup R_{\mathrm{blue}})\bigr)
          \cup X_{\mathrm{fill}}\cup Y_{\mathrm{blue}}.
\]
These are set unions.  Since the removal sets are disjoint and have
total size $r$, and the additions have total size at most $a+c=m$,
\begin{equation}
 |U'|\le h-|Q|-r+a+c=h.
 \label{comp:cardinality}
\end{equation}
This inequality permits repeated ball choices, overlapping balls,
coincidences with retained facilities, and reopening any member of $Q$
or either removal set.  Pad $U'$ to size $h$, if necessary, obtaining
$U_X$.  Padding is possible because $h\le|F|$ and makes $U_X$ nonempty.
Only now invoke domination of the permanent dummy: any real center is
within the original diameter of every original point.  A filled center
in $\mathcal F_i$ similarly satisfies
$d(y,X_{\mathrm{fill}})\le d(y,w_i)+R_i=d(y,\delta_i)$ for any original
point $y$.  Every point of $Z$ therefore has its distance from original
points attained or improved by the final solution.

For a client whose minimum in the base term uses $X$, the distance is
already attained by $U_X$.  Otherwise use its fixed assignment in $H$.
A retained real center or a dummy is covered by the preceding argument.
If its assigned center $s$ is removed red, triangle inequality bounds
its new distance by $d(p,H)+d(s,Z\cup X)$.  If it is removed blue, use
$d(p,H)+d(s,y_s)$ for its stored destination $y_s$.  Charging these
additional distances to all $w'_s$ clients of a removed center is an
upper bound even when some of their base minima use $X$.  Summing proves
$\operatorname{cost}(U_X)\le\Psi(X)$.

For submodularity, each client's base-cost saving is
\[
 d(p,H)-\min\{d(p,H),d(p,X)\}
 =\max\bigl(0,\max_{x\in X}\{d(p,H)-d(p,x)\}\bigr).
\]
The maximum of fixed offers has decreasing marginal gains and is
monotone.  If there are no red candidates the red term is zero.
Otherwise choose a common
$C\ge\max_{s\text{ red}}w'_sd(s,Z)$, and write
\[
 v_s(X)=C-w'_sd(s,Z\cup X)
       =\max\bigl(C-w'_sd(s,Z),
                    \max_{x\in X}\{C-w'_sd(s,x)\}\bigr).
\]
The scores $v_s(X)$ lie in $[0,C]$.  The sum of the
$q=b_{\mathrm{red}}$ smallest red costs equals $qC$ minus the sum
$V_q(X)$ of the $q$ largest scores.  If $C=0$ or $q=0$ this term is
constant.  Otherwise, counting scores above each level gives
\[
 V_q(X)=\int_0^C\min\bigl(q,|\{s:v_s(X)\ge t\}|\bigr)\,dt.
\]
At a fixed level $t$, each selected $x$ covers a fixed set of red
centers, in addition to a fixed baseline covered set.  The marginal of
adding $x$ to its coverage count truncated at $q$ is the minimum of
the number of newly covered centers and the remaining capacity below
$q$.  Both numbers can only decrease when the previously selected set
grows.  The integrand is therefore monotone and submodular, and so is
its integral.  The red saving is $V_q(X)-V_q(\varnothing)$.
Adding the base savings establishes the assertion.  Identical offers
from copies in different balls preserve this reasoning.
\end{proof}

\subsection{A completion inequality with arbitrary slack}

We have constructed a computable surrogate, proved that its savings
are submodular, and realized every feasible choice without exceeding
the facility budget.  We now analyze one favorable color-and-label
event.  The following statement retains the exact base contribution
outside the intended removals, which is essential to the refined
recovery argument.

\begin{theorem}[Completion of a prefix]
\label{comp:completion}
In the prefix setting above, fix a rational $1/e<\theta<1$.  Define
\[
 R^*=S_0\setminus Q,\qquad
 J=\bigcup_{s\in R^*}J_s,\qquad
 I=\sum_{p\in J}d(p,H),\qquad
 X^*=\{(i,i):i\in A\}.
\]
For $p\in C_i$, let
\begin{equation}
 c_p=\begin{cases}
 d(p,\phi_i),&i\text{ good},\\
 d_p^*+\theta(r_i+R_i),&i\in A,\\
 d_p^*,&i\in B,
 \end{cases}
 \qquad
 b_p=\theta d(p,H)+(1-\theta)\min\{d(p,H),d(p,X^*)\}.
 \label{comp:client-charges}
\end{equation}
There is an algorithm, using only $S,m,Q$ and the given balls, whose
outputs always have at most $h$ real facilities and whose work is
$\exp(O(m))\operatorname{poly}(N,1/(\theta-1/e))$, with polynomial
dependence also on the input bit length.  There is an event in its
color and label choices of probability at least $\exp(-O(m))$ such
that, for every realization of those choices in the event, its output
$U$ satisfies
\begin{equation}
 \mathbb E[\operatorname{cost}(U)]
       \le 2I+\sum_{p\in J}c_p+\sum_{p\notin J}b_p.
 \label{comp:completion-bound}
\end{equation}
The expectation is over the remaining optimization randomness.  For
good clients $b_p\le c_p$; for clients of $A$ also $b_p\le c_p$; for
clients of $B$ one always has $b_p\le d(p,H)$.  No upper bound on
$\operatorname{cost}(H)$, proxy error, or ball radii is needed for this
theorem beyond containment of $i$ in its designated ball.
\end{theorem}

\begin{proof}
For each feasible count branch, apply Lemma~\ref{comp:partition-max}
to the function in Lemma~\ref{comp:surrogate-properties}, with error
at most $\theta-1/e$, and construct the corresponding $U_X$.  Keep
the cheapest constructed solution, including $S$ as a fallback.  For
the fixed branch under consideration, the gain guarantee and the
feasibility of $X^*$ imply
\begin{equation}
 \mathbb E[\operatorname{cost}(U)]
 \le\mathbb E[\Psi(X)]
 \le\theta\Psi(\varnothing)+(1-\theta)\Psi(X^*).
 \label{comp:convex-surrogate}
\end{equation}
It remains to identify a suitable branch and bound this convex
combination.

For each $s\in R^*$, minimize the following mixed distance over all
comparison destinations, resolving ties by a fixed rule:
\begin{equation}
 M_s=\min\left\{
 \begin{array}{ll}
 d(s,\phi_i),&i\text{ good},\\
 (1-\theta)d(s,i)+\theta\bigl(d(s,w_i)+R_i\bigr),&i\in A,\\
 d(s,i),&i\in B.
 \end{array}\right.
 \label{comp:mixed-distance}
\end{equation}
Require $s$ to be red for either of the first two destination types and
blue for the third.  Require every proxy used in the first type to be
green.  The requirements are consistent: $R^*\subseteq S_0$ whereas
all proxies lie outside $S_0$.  At most $r$ removal centers and at most
$r$ proxy destinations receive specified colors, so the event has
probability at least $4^{-2m}$.  All other proxies may have any color;
if the optimization removes one of them, its displacement is paid by
the surrogate just as for any other removed center.

Use the count branch equal to the sizes of these designated red and
blue sets.  If the designated blue count is $b>0$, require their labels
to be distinct.  Its conditional probability is at least
\[
 b!(2b)^{-b}\ge(2e)^{-b}.
\]
The inequality follows, for example, from
$\log(b!)\ge\int_1^b\log x\,dx=b\log b-b+1$.
For $b=0$ the event has probability one.  The blue comparison assigns
each designated blue removal to its chosen member of $B$.  Grouping
equal destinations uses at most $|B|=c$ blocks, so it is feasible for
Equation~\eqref{comp:blue-block} and the dynamic program.  Consequently
its $K$ is at most the sum of designated blue charges $w'_sM_s$.
In particular the successful branch is never infeasible.

For the designated red subset, at $X=\varnothing$ a proxy destination
is green and a ball destination has its leaf in $Z$.  At $X=X^*$ the
corresponding true ball center is also available.  At each of these
two arguments, the optimized sum of smallest red costs is no greater
than the sum over this same designated red subset.  Therefore their
convex combination, together with $K$, is at most
\begin{equation}
 \sum_{s\in R^*}w'_sM_s.
 \label{comp:removal-comparison}
\end{equation}
The minimizing red subset need not be the same at the two arguments;
the fixed designated subset is an admissible upper bound at both.

For $p\in J_s\cap C_i$ offer its own cluster as a destination in
Equation~\eqref{comp:mixed-distance}.  In the good and uncovered cases,
triangle inequality gives respectively
$M_s\le d(p,H)+d(p,\phi_i)$ and
$M_s\le d(p,H)+d_p^*$.  In the covered case it gives
\[
 \begin{split}
 M_s&\le(1-\theta)d(s,i)+\theta(d(s,w_i)+R_i)\\
 &\le d(p,H)+(1-\theta)d_p^*
            +\theta(d_p^*+r_i+R_i)
   =d(p,H)+c_p.
 \end{split}
\]
Although $M_s$ is chosen once for all clients of $J_s$, it is no larger
than every offered value.  Summing therefore bounds
Equation~\eqref{comp:removal-comparison} by $I+\sum_{p\in J}c_p$.
The base terms of the convex combination contribute at most a further
$I$ on $J$, and contribute exactly $b_p$ on its complement.  Together
with Equation~\eqref{comp:convex-surrogate} this proves
Equation~\eqref{comp:completion-bound}.

For a good client $d(p,H)\le d(p,\phi_i)$ since $Q\subseteq S_0$.
For a covered client, $d(p,H)\le d_p^*+r_i+R_i$ and
$d(p,X^*)\le d_p^*$, yielding $b_p\le c_p$.  The bound for $B$ follows
from the definition of $b_p$.  All complexity statements follow from
the $r+1$ count branches, the subset dynamic program, and
Lemma~\ref{comp:partition-max}.  The number of copied facilities is at
most $a|F|\le N^2$, and the surrogate can be evaluated using distance
lookups and sorting.  The color and label constructions use bounded
numbers of fair bits.  Combining their probabilities gives the claimed
$\exp(-O(m))$ event.
\end{proof}

\subsection{A quantitative recovery guarantee}

\begin{theorem}[Single-exponential recovery]
\label{comp:recovery}
Suppose $|S|=|O|=h\ge1$, $P=\operatorname{cost}(O)>0$, and $O$ has
an injective proxy map into $S$ for its good centers such that
\[
 \sum_{i\text{ good}}\sum_{p\in C_i}d(p,\phi_i)
       \le P_g+0.01P.
\]
Assume the number $m$ of bad centers is at most a supplied bound
$M_0$.  On the positive bounded integral instances used in the body
of the proof, for every $0<\zeta<1$ an algorithm knowing only
$S,M_0,\zeta$ returns a real solution of size at most $h$ and cost at
most $1.97P$ with probability at least $1-\zeta$.  Its running time is
\[
 \exp(O(M_0))\operatorname{poly}(N)\log(2/\zeta).
\]
In particular this is polynomial time for every fixed constant in
$M_0=O(\log N)$.  Local optimality of $S$ is not an additional
hypothesis of this theorem.
\end{theorem}

\begin{proof}
If $m=0$, the cost of $S$ is at most the displayed proxy-assignment
cost, hence at most $1.01P$.  The algorithm always retains $S$ as a
candidate.  Suppose henceforth $m>0$, and fix
$\kappa=10^{-6}$ and $\theta=0.38$.  Run the sampling procedure of
Theorem~\ref{sam:successful-prefix} for at most $2m$ steps, examining
every prefix.  Its successful prefix has the structure used in
Theorem~\ref{comp:completion}, and satisfies
\[
 I\le\kappa P,\qquad
 T_i:=\sum_{p\in N_i}d(p,H),\quad
 \sum_{i\in B}T_i\le\kappa P,
 \qquad N_i=\{p\in C_i:d_p^*\le1.2\bar d_i\}.
\]
Here $\bar d_i=P_i/n_i$ for $n_i>0$; empty clusters contribute no
clients and no terms with denominator $n_i$.  The same theorem gives
\begin{equation}
 \begin{gathered}
 r_i\le1.2\bar d_i,\qquad
 R_i\le(1+\kappa)r_i+\kappa r^0(w_i),\qquad
 \sum_{i\in A}n_i r^0(w_i)\le5.2P,\\
 d(p,H)\le1.85d_p^*+60T_i/n_i
       \quad(i\in B,\ p\in C_i\setminus N_i).
 \end{gathered}
 \label{comp:main-endpoint}
\end{equation}

Apply Theorem~\ref{comp:completion}.  Good clients, whether inside or
outside $J$, contribute at most their proxy-assignment cost.  Covered
bad clients contribute at most
\[
 \sum_{i\in A}\left[P_i+
   \theta n_i\bigl((2+\kappa)1.2\bar d_i+\kappa r^0(w_i)\bigr)\right].
\]
For a remaining bad cluster, clients in $J$ cost at most $d_p^*$;
clients outside $J$ and inside $N_i$ contribute at most $d(p,H)$,
whose sum is bounded by $T_i$; the other clients are bounded by the
last line of Equation~\eqref{comp:main-endpoint}.  Thus their total is
at most $1.85\sum_{i\in B}P_i+61\sum_{i\in B}T_i$.
Including the two copies of $I$ gives the conditional expectation
\[
 \begin{split}
 \mathbb E[\operatorname{cost}(U)]
 &\le 2\kappa P+P_g+0.01P\\
 &\quad+\sum_{i\in A}\left[P_i+
   \theta n_i\bigl((2+\kappa)1.2\bar d_i+\kappa r^0(w_i)\bigr)\right]
       +1.85\sum_{i\in B}P_i+61\kappa P\\
 &\le\left[1+0.38(2+10^{-6})1.2+0.01
                 +63\cdot10^{-6}+0.38\cdot5.2\cdot10^{-6}\right]P\\
 &=1.922065432P<1.94P.
 \end{split}
\]
In the second inequality the covered-cluster coefficient
$1+0.38(2+\kappa)1.2=1.912000456$ bounds both the good coefficient
$1$ and the uncovered coefficient $1.85$.

The output cost is nonnegative.  Markov's inequality therefore gives,
conditional on the successful prefix and the color and label event,
\[
 \Pr\{\operatorname{cost}(U)\le1.97P\}
 \ge 1-\frac{1.94}{1.97}=\frac3{197}.
\]
The sampling event has probability $\exp(-O_\kappa(m))$ by
Theorem~\ref{sam:successful-prefix}; the completion event has probability
$\exp(-O(m))$.  Their product with $3/197$ is at least
$\exp(-C m)$ for an absolute constant $C$ after enlarging $C$, since
$\kappa$ is fixed and $m\ge1$.  These are conditional bounds for every
successful prefix, so they combine without an independence assumption
between its shape and the required colors.

The algorithm need not know the comparison solution, its partition,
its proxies, or $P$.  It enumerates $1\le m\le\min\{M_0,h\}$ and
all integer cost guesses $1\le P\le n\Delta$, where $\Delta$ is the
maximum original distance of the bounded integral instance.  One guess
equals the true $P$.  For each guess repeat the full sampling and
completion experiment
$\lceil\exp(C M_0)\log(1/\zeta)\rceil$ times, taking the best feasible
candidate over all prefixes and all guesses, including $S$.
The chance that every trial for the correct guess fails is at most
$\zeta$.  The number of cost guesses is polynomial in $N$, each
sampling trial has polynomial work, and completion contributes only
$\exp(O(m))\operatorname{poly}(N)$.  The stated running time follows.
All branches, including incorrect guesses, discard prefixes with
$|Q|>m$ or more than $m$ balls, empty facility balls, or impossible
counts, and obey Equation~\eqref{comp:cardinality};
the fallback also has size $h$.  Consequently the budget holds on
every outcome.  The algorithm retains the minimum actual cost without
testing any condition involving $O$ or recognizing successful events.
\end{proof}
